\section*{Bab \ref{ch:g12:light-as-wave} --- Cahaya sebagai Gelombang: Difraksi dan Interferensi}

\begin{solution}{exo:g12:light-as-wave:1}
$f = c/\lambda$: merah $\num{3.00e8}/\num{7.0e-7} \approx
\qty{4.3e14}{Hz}$; ungu $\approx \qty{7.5e14}{Hz}$ --- ungulah yang
lebih besar (\omterm{def:g12:mechanical-waves:wavelength}{panjang gelombangnya} lebih pendek, lajunya sama, yaitu $c$).
\end{solution}

\begin{solution}{exo:g12:light-as-wave:2}
$\theta = \num{5.32e-7}/\num{2.0e-4} = \qty{2.7e-3}{rad}$;
$L = 2\theta D = 2 \times \num{2.66e-3} \times 1.5 \approx
\qty{8.0}{mm}$.
\end{solution}

\begin{solution}{exo:g12:light-as-wave:3}
Suara: $\lambda = 340/425 = \qty{0.80}{m}$, sehingga $\lambda/a = 1.0$
--- \omterm{def:g12:light-as-wave:diffraction}{difraksinya} terbesar, dan \omterm{def:g10:signals-and-waves:sound}{bunyinya} membanjiri lorongnya. Cahaya:
$\lambda/a \approx \num{5.5e-7}/0.80 \approx \num{7e-7}$ ---
penyebarannya dapat diabaikan: cahayanya tetap lurus, sehingga sudut
ruangannya menyembunyikannya.
\end{solution}

\begin{solution}{exo:g12:light-as-wave:4}
Di $M$: $\delta = \qty{6.0}{cm} = 3\lambda$, \omterm{prop:g12:light-as-wave:conditions}{konstruktif} --- nyaring.
Di $N$: $\delta = \qty{5.0}{cm} = (2 + \tfrac12)\lambda$, \omterm{prop:g12:light-as-wave:conditions}{destruktif} ---
(nyaris) senyap.
\end{solution}

\begin{solution}{exo:g12:light-as-wave:5}
$i = \lambda D/a = \num{6.33e-7} \times 1.8/\num{2.5e-4} \approx
\qty{4.6}{mm}$. Di dalam $\pm\qty{1.0}{cm}$: $|k| \leq 10/4.6 = 2.2$,
sehingga $k = -2, \dots, 2$: ada lima pita terang.
\end{solution}

\begin{solution}{exo:g12:light-as-wave:6}
$d = 2\lambda D/L = 2 \times \num{6.5e-7} \times 1.6/0.028 \approx
\qty{74}{\micro m}$.
\end{solution}

\begin{solution}{exo:g12:light-as-wave:7}
$i = \qty{4.7}{mm}$, sehingga $\lambda = i\,a/D = \num{4.7e-3} \times
\num{2.0e-4}/1.5 \approx \qty{6.3e-7}{m} = \qty{630}{nm}$: merah.
Mengukur sepuluh \omterm{def:g12:light-as-wave:fringes}{jarak antarpita} membagi galat pembacaan penggarisnya
dengan sepuluh.
\end{solution}

\begin{solution}{exo:g12:light-as-wave:8}
Kedua lampunya (secara kasar) memiliki \omterm{def:g10:signals-and-waves:frequency}{frekuensi} yang sepadan tetapi
\emph{tidak seirama}: masing-masing memancarkan rangkaian \omterm{def:g10:signals-and-waves:wave}{gelombang} yang
pendek dengan fase yang melompat secara acak, sehingga sumbernya tidak
\omterm{def:g12:light-as-wave:coherent}{koheren} dan pita sesaatnya bergeser miliaran kali tiap sekon --- mata
melihat rata-ratanya, yang seragam. Young menyinari kedua celahnya dari
\emph{satu} \omterm{def:g10:signals-and-waves:wave}{gelombang}: kedua salinannya mewarisi setiap lompatan fasenya
bersama-sama lalu tetap seirama.
\end{solution}

\begin{solution}{exo:g12:light-as-wave:9}
Dari yang pertama sampai yang ketujuh $= 6i = \qty{18.0}{mm}$, sehingga
$i = \qty{3.0}{mm}$;
$\lambda = i\,a/D = \num{3.0e-3} \times \num{3.0e-4}/1.4 \approx
\qty{6.4e-7}{m} \approx \qty{640}{nm}$.
\end{solution}

\begin{solution}{exo:g12:light-as-wave:10}
$\sin\theta = k\lambda/d = 0.396\,k$: untuk $k = 1$,
$\theta = 23^\circ$; untuk $k = 2$, $\sin\theta = 0.79$,
$\theta = 52^\circ$; $k = 3$ akan menuntut $\sin\theta = 1.19 > 1$:
mustahil. Orde tertingginya $k = 2$.
\end{solution}

\begin{solution}{exo:g12:light-as-wave:11}
$\theta = \num{5.5e-7}/0.10 = \qty{5.5e-6}{rad}$. Kedua lampu depannya
membentang $1.5/\num{1.0e5} = \qty{1.5e-5}{rad}$, sekitar tiga kali
kaburannya: terpisah, walaupun pas-pasan.
\end{solution}

\begin{solution}{exo:g12:light-as-wave:12}
Pantulan dari muka depan dan muka belakangnya \omterm{prop:g12:light-as-wave:conditions}{berinterferensi}; di dalam
\omterm{def:g10:light-spectra:white}{cahaya putih}, tiap tebalnya menguatkan sebagian \omterm{def:g12:mechanical-waves:wavelength}{panjang gelombang} dan
menghapus yang lain, sehingga muncullah sebuah warna. Tebal selaputnya
hanya bergantung pada ketinggiannya, sehingga titik yang warnanya sama
membentuk pita mendatar. Penirisannya menipiskan selaputnya, sehingga
tebal yang melukis warna tertentu berada makin ke bawah: pitanya pun
merayap turun.
\end{solution}

\begin{solution}{exo:g12:light-as-wave:13}
$i = \lambda D/a$: ungu $\num{4.0e-7} \times 1.5/\num{2.0e-4} =
\qty{3.0}{mm}$; merah $\qty{5.6}{mm}$. Di tengahnya: semua warnanya
terang --- satu pita putih. Beberapa milimeter ke luar, pola warnanya
sudah saling menjauh: pita yang berwarna pelangi. Lebih jauh lagi,
puncak semua warnanya bertindihan di mana-mana: kaburan putih yang
seragam.
\end{solution}

\begin{solution}{exo:g12:light-as-wave:14}
$\theta = \num{5.5e-7}/2.4 = \qty{2.3e-7}{rad}$; pada \qty{250}{km}
nilai itu berarti $\num{2.3e-7} \times \num{2.5e5} \approx \qty{6}{cm}$.
\omterm{def:g10:universal-gravitation:satellite}{Satelit} itu dapat menghitung mobil dan menemukan seseorang, tetapi
huruf cetak koran (berukuran milimeter) dua puluh kali di bawah batas
\omterm{def:g12:light-as-wave:diffraction}{difraksinya}: pengakuannya hanya mitos.
\end{solution}

\begin{solution}{exo:g12:light-as-wave:15}
\omterm{def:g10:signals-and-waves:frequency}{Frekuensinya} ditetapkan sumbernya: tidak berubah. Lajunya turun menjadi
$c/n$, sehingga $\lambda' = \lambda/n$ dan
$i' = i/n = 4.0/1.33 \approx \qty{3.0}{mm}$: pitanya merapat.
\end{solution}

\begin{solution}{pb:g12:light-as-wave:1}
\textbf{1.} $\theta = \lambda/a = \num{6.33e-7}/\num{1.00e-4} =
\qty{6.3e-3}{rad}$.

\textbf{2.} Setengah lebarnya $D\tan\theta \approx D\theta = \lambda D/a$,
sehingga $L = 2\lambda D/a = 2 \times \num{6.33e-7} \times
3.00/\num{1.00e-4} = \qty{3.8}{cm}$.

\textbf{3.} $(3.8 - 3.7)/3.8 \approx 3\%$: masih di dalam ketelitian
sebuah penggaris --- caranya terbukti.

\textbf{4.} Karena $L \propto 1/a$: menyeparuhkan $a$ menggandakan $L$,
sehingga $L = \qty{7.6}{cm}$. Makin sempit celah atau penghalangnya,
makin lebar polanya.

\textbf{5.} $\tan(\num{6.33e-3}) = \num{6.33008e-3}$: keduanya sesuai
sampai sekitar $\num{1e-5}$ secara relatif --- hampirannya jauh lebih
baik daripada pengukuran mana pun di sini.

\textbf{6.} Menurut \cref{prop:g12:light-as-wave:hair}: sebuah
penghalang selebar $d$ mendifraksikan seperti celah selebar $d$,
sehingga $L = 2\lambda D/d$ tetap berlaku.

\textbf{7.} $d = 2\lambda D/L = 2 \times \num{6.33e-7} \times
3.00/0.054 \approx \qty{70}{\micro m}$.

\textbf{8.} $L = \qty{5.2}{cm} \to d = \qty{73}{\micro m}$;
$L = \qty{5.6}{cm} \to d = \qty{68}{\micro m}$: sehingga
$d = \qty{70 \pm 3}{\micro m}$.

\textbf{9.} $d = \num{3.80e-6}/0.076 = \qty{50}{\micro m}$: rambut
temanyalah yang lebih halus --- dan karena $L \propto 1/d$, rambut yang
lebih halus melemparkan pola yang lebih lebar.

\textbf{10.} Tidak: sebuah lampu berwarna putih (tiap \omterm{def:g12:mechanical-waves:wavelength}{panjang
gelombangnya} melukis pola yang berbeda, lalu semuanya kabur) sekaligus
tidak \omterm{def:g12:light-as-wave:coherent}{koheren} dan lebar (tidak ada satu berkas bersih untuk
didifraksikan). \omterm{def:g11:color-light-sources:lamps}{Lasernya} \omterm{def:g10:light-spectra:white}{monokromatik} dan terarah.

\textbf{11.} Hanya salinan satu \omterm{def:g10:signals-and-waves:wave}{gelombang} yang \omterm{def:g12:light-as-wave:coherent}{koheren}; dua sumber yang
saling bebas hanyut keluar dari iramanya dan pitanya pun luntur
(\cref{rem:g12:light-as-wave:lamps}).

\textbf{12.} Di tengahnya $\delta = 0 = 0 \times \lambda$ untuk setiap
$\lambda$: \omterm{prop:g12:light-as-wave:conditions}{konstruktif} tanpa menghiraukan \omterm{def:g12:mechanical-waves:wavelength}{panjang gelombangnya}.

\textbf{13.} $i = 4.26/10 = \qty{4.3}{mm}$; satu rentang berisi sepuluh
dibaca dengan galat penggaris yang sama seperti satu \omterm{def:g12:light-as-wave:fringes}{jarak antarpita},
sehingga galat pada $i$ terbagi sepuluh.

\textbf{14.} $\lambda = i\,a'/D = \num{4.26e-3} \times
\num{2.50e-4}/2.00 = \num{5.33e-7} \approx \qty{533}{nm}$.

\textbf{15.} Nilai $\qty{533}{nm}$ terletak di dalam
$\qty{532 \pm 10}{nm}$: sesuai; \omterm{def:g12:mechanical-waves:wavelength}{panjang gelombang} itu berwarna hijau.

\textbf{16.} $\tan\theta = 0.43/1.00$, sehingga $\theta = 23^\circ$. Di
sini $\tan\theta = 0.43$ sedangkan $\sin\theta = 0.40$: pada
$23^\circ$, penyamaan $\sin$, $\tan$ dan $\theta$ untuk sudut kecil
meleset hampir $10\%$ --- pakailah fungsinya yang tepat.

\textbf{17.} $d = \lambda/\sin\theta = \num{6.33e-7}/0.396 \approx
\qty{1.6}{\micro m}$.

\textbf{18.} $\tan\theta = 1.65$, $\theta = 59^\circ$,
$\sin\theta = 0.855$: $d = \num{6.33e-7}/0.855 \approx
\qty{0.74}{\micro m}$.

\textbf{19.} $1.6/0.74 \approx 2.2$: alur DVD dua kali lebih rapat, dan
lubangnya pun sepadan lebih kecil (dibaca dengan \omterm{def:g11:color-light-sources:lamps}{laser} yang \omterm{def:g12:mechanical-waves:wavelength}{panjang
gelombangnya} lebih pendek) --- beberapa kali lipat datanya di dalam
cakram \qty{12}{cm} yang sama.

\textbf{20.} Rambutnya \qty{70 \pm 3}{\micro m}; rambut temannya
\qty{50}{\micro m}; \omterm{def:g11:color-light-sources:lamps}{laser} hijaunya \qty{533}{nm}; alur CD-nya
\qty{1.6}{\micro m}; alur DVD-nya \qty{0.74}{\micro m} --- satu \omterm{def:g12:mechanical-waves:wavelength}{panjang
gelombang} yang diketahui mengubah ilmu ukur setiap polanya menjadi
sebuah panjang, sehingga cahaya setengah mikrometer menjadi penggaris
yang berskala setengah mikrometer.
\end{solution}
